% ==============================================================================
%  CHAPTER 4 — RESULTS
% ==============================================================================
\section{Results}
\label{sec:results}

This chapter reports the results on the three labelled sets of Table~\ref{tab:gtsets},
the effect of the offline second pass, and the processing speed. Interpretation is left
to Chapter~\ref{sec:discussion}.

\subsection{Phase 2a: Geometry Cues Only}
\label{sec:results_2a}

With only the proximity and centrality cues, the system was evaluated on the 66 tracks of
Set~A (Table~\ref{tab:confusion_2a}). It blurred 35 of 40 bystanders and kept 22 of 26
partners visible.

\begin{table}[htbp]
\begin{singlespace}
\centering
\caption{Phase~2a confusion matrix on Set~A. Rows give the human label and columns the
system's decision.}
\label{tab:confusion_2a}
\renewcommand{\arraystretch}{1.3}
\begin{tabular}{@{}lccc@{}}
\toprule
 & \textbf{Kept visible} & \textbf{Blurred} & \textbf{Total} \\
\midrule
\textbf{Partner}   & 22 (TP) & 4 (FN)  & 26 \\
\textbf{Bystander} & 5 (FP)  & 35 (TN) & 40 \\
\textbf{Total}     & 27      & 39      & 66 \\
\bottomrule
\end{tabular}
\end{singlespace}
\end{table}

This gives PSR = 87.50\% (95\% CI 73.9-94.5\%), IPR = 84.62\% (66.5-93.9\%),
precision 81.48\% and F$_1$ 83.02\%.

\subsection{Phase 2b: Still-Frame Evaluation}
\label{sec:results_2b}

With the head-pose and speaking cues added, the system was first evaluated on the 29 faces
of Set~B. The four annotators gave 114 of the 116 possible labels; the majority vote
produced 7 partners and 22 bystanders, with one tie. Pairwise Cohen's~$\kappa$ ranged
from 0.307 to 0.724 (mean 0.587), so agreement was moderate.

\begin{table}[htbp]
\begin{singlespace}
\centering
\caption{Phase~2b confusion matrix on Set~B (still frames).}
\label{tab:confusion_2b}
\renewcommand{\arraystretch}{1.3}
\begin{tabular}{@{}lccc@{}}
\toprule
 & \textbf{Kept visible} & \textbf{Blurred} & \textbf{Total} \\
\midrule
\textbf{Partner}   & 4 (TP) & 3 (FN)  & 7 \\
\textbf{Bystander} & 0 (FP) & 22 (TN) & 22 \\
\textbf{Total}     & 4      & 25      & 29 \\
\bottomrule
\end{tabular}
\end{singlespace}
\end{table}

All 22 bystanders were blurred, including the 8 who appeared in interaction videos next
to a partner, giving PSR = 100\% (85.1-100\%). Four of the 7 partners were kept visible,
giving IPR = 57.14\% (25.0-84.2\%), precision 100\% and F$_1$ 72.73\%. With only seven
partners, each one changes IPR by 14 percentage points, which is why its interval is so
wide.

\subsection{Phase 2b: Video-Clip Evaluation}
\label{sec:results_clips}

Set~C repeats the Phase~2b evaluation on 102 clips with sound. Agreement between the five
annotators was almost perfect: pairwise Cohen's~$\kappa$ ranged from 0.818 to 0.957, and
Fleiss'~$\kappa$ was 0.887 (Table~\ref{tab:clip_kappa}). No clip ended in a tie.

\begin{table}[htbp]
\begin{singlespace}
\centering
\caption{Annotator agreement on Set~C. Annotator A is myself.}
\label{tab:clip_kappa}
\renewcommand{\arraystretch}{1.3}
\begin{tabular}{@{}lcc@{}}
\toprule
\textbf{Annotator pair} & \textbf{Cohen's $\kappa$} & \textbf{Raw agreement} \\
\midrule
A and B & 0.957 & 98.0\% \\
A and C & 0.913 & 96.1\% \\
A and D & 0.818 & 92.1\% \\
A and E & 0.822 & 92.2\% \\
B and C & 0.956 & 98.0\% \\
B and D & 0.861 & 94.1\% \\
B and E & 0.864 & 94.1\% \\
C and D & 0.905 & 96.0\% \\
C and E & 0.907 & 96.1\% \\
D and E & 0.852 & 94.1\% \\
\midrule
\textbf{All five (Fleiss' $\kappa$)} & \textbf{0.887} & \\
\bottomrule
\end{tabular}
\end{singlespace}
\end{table}

Tables~\ref{tab:confusion_clips} and~\ref{tab:metrics_clips} give the results.

\begin{table}[htbp]
\begin{singlespace}
\centering
\caption{Phase~2b confusion matrix on Set~C (video clips).}
\label{tab:confusion_clips}
\renewcommand{\arraystretch}{1.3}
\begin{tabular}{@{}lccc@{}}
\toprule
 & \textbf{Kept visible} & \textbf{Blurred} & \textbf{Total} \\
\midrule
\textbf{Partner}   & 39 (TP) & 30 (FN) & 69 \\
\textbf{Bystander} & 1 (FP)  & 32 (TN) & 33 \\
\textbf{Total}     & 40      & 62      & 102 \\
\bottomrule
\end{tabular}
\end{singlespace}
\end{table}

\begin{table}[htbp]
\begin{singlespace}
\centering
\caption{Phase~2b metrics on Set~C with 95\% Wilson intervals. IPR, precision and F$_1$
are not defined for bystander videos, which contain no partners.}
\label{tab:metrics_clips}
\renewcommand{\arraystretch}{1.3}
\fitwidth{\begin{tabular}{@{}lccc@{}}
\toprule
\textbf{Metric} & \textbf{All clips} & \textbf{Interaction videos} & \textbf{Bystander videos} \\
\midrule
Faces evaluated & 102 & 81 & 21 \\
PSR       & 96.97\% (84.7-99.5) & 100.00\% (75.7-100.0) & 95.24\% (77.3-99.2) \\
IPR       & 56.52\% (44.8-67.6) & 56.52\% (44.8-67.6) & \\
Precision & 97.50\% & 100.00\% & \\
F$_1$     & 71.56\% & 72.22\% & \\
\bottomrule
\end{tabular}}
\end{singlespace}
\end{table}

The system blurred 32 of 33 bystanders, including all 12 who appeared in interaction
videos alongside the wearer's group, and kept 39 of 69 partners visible. Of the 40 faces
it kept visible, 39 were partners. Among the 30 partners it blurred, 25 had a median face
size below the 1.2\% area threshold, compared with a median of 1.72\% for the partners it
kept visible, and 10 were confirmed as partners later in the same video, after their clip
had ended.

\subsection{Offline Second Pass}
\label{sec:results_secondpass}

Table~\ref{tab:secondpass} shows the effect of the two rules of
Section~\ref{sec:secondpass} on Set~C.

\begin{table}[htbp]
\begin{singlespace}
\centering
\caption{Effect of the offline second pass on Set~C (102 clips, 95\% Wilson intervals).}
\label{tab:secondpass}
\renewcommand{\arraystretch}{1.3}
\fitwidth{\begin{tabular}{@{}lccc@{}}
\toprule
\textbf{Metric} & \textbf{Live system} & \textbf{+ Retroactive latch} & \textbf{+ Group rule} \\
\midrule
TP / FN / FP / TN & 39 / 30 / 1 / 32 & 49 / 20 / 2 / 31 & 54 / 15 / 2 / 31 \\
PSR       & 96.97\% (84.7-99.5) & 93.94\% (80.4-98.3) & 93.94\% (80.4-98.3) \\
IPR       & 56.52\% (44.8-67.6) & 71.01\% (59.4-80.4) & 78.26\% (67.2-86.4) \\
Precision & 97.50\% & 96.08\% & 96.43\% \\
F$_1$     & 71.56\% & 81.67\% & 86.40\% \\
\bottomrule
\end{tabular}}
\end{singlespace}
\end{table}

The retroactive latch recovered 10 partners and exposed one more bystander. The group rule
recovered 5 further partners with no further exposure. I repeated the evaluation with the
group rule's time limit set anywhere from 10 to 60~seconds: PSR stayed at 93.94\% and IPR
ranged from 72.46\% to 78.26\%. Only a limit as short as 5~seconds admitted another
bystander.

\subsection{Comparison with the Baseline}
\label{sec:baseline_comparison}

The blanket-blurring baseline blurs every face, so by construction its PSR is 100\% and its
IPR is 0\%. Table~\ref{tab:vs_baseline} places each configuration beside it.

\begin{table}[htbp]
\begin{singlespace}
\centering
\caption{All configurations against the blanket-blurring baseline.}
\label{tab:vs_baseline}
\renewcommand{\arraystretch}{1.3}
\begin{tabular}{@{}llccc@{}}
\toprule
\textbf{Configuration} & \textbf{Set} & \textbf{PSR} & \textbf{IPR} & \textbf{F$_1$} \\
\midrule
Blanket blurring (baseline)      & any & 100\%    & 0\%     & 0\% \\
Phase~2a, geometry cues          & A   & 87.50\%  & 84.62\% & 83.02\% \\
Phase~2b, all cues               & B   & 100.00\% & 57.14\% & 72.73\% \\
Phase~2b, all cues               & C   & 96.97\%  & 56.52\% & 71.56\% \\
Phase~2b with offline second pass & C  & 93.94\%  & 78.26\% & 86.40\% \\
\bottomrule
\end{tabular}
\end{singlespace}
\end{table}

\subsection{Processing Speed}
\label{sec:fps}

Table~\ref{tab:fps} gives the time spent in each stage of the Phase~2a pipeline on the
laptop CPU under mains power.

\begin{table}[htbp]
\begin{singlespace}
\centering
\caption{Mean time per frame for each stage (Phase~2a, 300 frames, laptop CPU, mains power).}
\label{tab:fps}
\renewcommand{\arraystretch}{1.3}
\begin{tabular}{@{}lcc@{}}
\toprule
\textbf{Stage} & \textbf{Time (ms)} & \textbf{Share of total} \\
\midrule
Frame reading                  & 10.0  & 11.6\% \\
Face detection (YOLOv8n-face)  & 55.1  & 64.1\% \\
Tracking (SORT)                & 20.1  & 23.4\% \\
Classification                 & 0.02  & $<$0.1\% \\
Selective blurring             & 0.87  & 1.0\% \\
\midrule
\textbf{Total per frame}       & \textbf{86.0} & \textbf{100\%} \\
\textbf{Throughput}            & \multicolumn{2}{c}{\textbf{11.62 FPS}} \\
\bottomrule
\end{tabular}
\end{singlespace}
\end{table}

Face detection takes almost two thirds of the time; the classifier itself takes 0.02~ms.
The landmark stage of Phase~2b, timed separately in two runs, took 20.1~ms per face crop
on average (21.1 and 19.1~ms). The test video had 0.99 tracked faces per frame, which adds
19.9~ms per frame and brings Phase~2b to about 106~ms per frame, or 9.4~FPS. The fastest
direct Phase~2b run reached 9.9~FPS. Because landmarks are computed for every tracked
face, this cost grows with the number of faces in view.

On battery power the laptop was slower: landmark extraction took 33.5~ms per face, Phase~2b
ran at 6.1~FPS and two Phase~2a runs gave 9.1 and 7.5~FPS. The processor also warmed up
over these consecutive runs, so these figures are a worst case rather than a precise
measurement.

% ── Power measurement: enable this paragraph after running eval/measure_power.py ──
% \paragraph{Power draw.}
% With the laptop on battery and all other applications closed, the idle laptop drew
% X.X~W. Running the Phase~2a pipeline raised this to X.X~W and Phase~2b to X.X~W, so the
% pipeline itself drew about X.X~W and X.X~W, or X.XX and X.XX joules per processed frame.

\paragraph{Detection on other hardware.}
Because detection dominates the processing time, Table~\ref{tab:det_hw} compares the
detection time I measured with published measurements of the same class of model on other
processors.

\begin{table}[htbp]
\begin{singlespace}
\centering
\caption{Time to detect faces in one $640 \times 640$ frame with a nano-sized YOLO model on
different processors. The laptop figure is my measurement; the others are published
benchmarks.}
\label{tab:det_hw}
\renewcommand{\arraystretch}{1.3}
\fitwidth{\begin{tabular}{@{}llllr@{}}
\toprule
\textbf{Processor} & \textbf{Unit} & \textbf{Format} & \textbf{Model} & \textbf{Time (ms)} \\
\midrule
Intel Core i5-13420H (laptop)        & CPU & PyTorch & YOLOv8n-face & 55.1 \\
Raspberry Pi 5 \cite{ultralytics_rpi} & CPU & PyTorch & YOLOv8n      & 508.6 \\
Raspberry Pi 5 \cite{ultralytics_rpi} & CPU & ONNX    & YOLOv8n      & 198.7 \\
Raspberry Pi 5 \cite{ultralytics_rpi} & CPU & NCNN    & YOLOv8n      & 94.3 \\
Snapdragon 8 Elite Gen 5 phone \cite{ultralytics_qnn} & CPU & QNN & YOLO26n & 53.3 \\
Snapdragon 8 Elite Gen 5 phone \cite{ultralytics_qnn} & GPU & QNN & YOLO26n & 17.2 \\
Snapdragon 8 Elite Gen 5 phone \cite{ultralytics_qnn} & NPU & QNN & YOLO26n & 11.3 \\
\bottomrule
\end{tabular}}
\end{singlespace}
\end{table}

\paragraph{Projected throughput on edge devices.}
Using the method of Section~\ref{sec:evaldesign}, Table~\ref{tab:projection} projects the
full pipeline to a Raspberry Pi~5. The detection time is taken from
Table~\ref{tab:det_hw}, and the remaining stages (31.0~ms for Phase~2a, plus 19.9~ms of
landmarks for Phase~2b) are multiplied by 3.5. For the Ray-Ban Meta glasses, whose
Snapdragon AR1 has four CPU cores at 1.9~GHz \cite{qcom_ar1}, I scaled the Raspberry Pi~5
projection by the ratio of clock speeds (2.4/1.9). This assumes the AR1's cores do about
as much work per clock cycle as the Pi's Cortex-A76 cores, which Qualcomm does not
publish. The last row places the detector on a smartphone-class neural processing unit
while the other stages stay on a Pi-class CPU; since the AR1's neural processing unit is
smaller than a flagship phone's, this row is an upper bound. The chips in Snap Spectacles
and Meta Aria Gen~2 are not published, so I could not project for them.

\begin{table}[htbp]
\begin{singlespace}
\centering
\caption{Measured and projected throughput of the full pipeline.}
\label{tab:projection}
\renewcommand{\arraystretch}{1.3}
\fitwidth{\begin{tabular}{@{}llcc@{}}
\toprule
\textbf{Device and configuration} & \textbf{Basis} & \textbf{Phase 2a} & \textbf{Phase 2b} \\
\midrule
Laptop, mains power                       & Measured  & 11.6 FPS & 9.4 FPS \\
Laptop, battery power                     & Measured  & 7.5-9.1 FPS & 6.1 FPS \\
Raspberry Pi~5, PyTorch detector          & Projected & 1.6 FPS  & 1.5 FPS \\
Raspberry Pi~5, NCNN detector             & Projected & 4.9 FPS  & 3.7 FPS \\
Ray-Ban Meta CPU, NCNN detector           & Projected & 3.9 FPS  & 2.9 FPS \\
Pi-class CPU with phone-class NPU detector & Upper bound & 8.3 FPS & 5.3 FPS \\
\bottomrule
\end{tabular}}
\end{singlespace}
\end{table}
